\emph{(o).} $\beta^{2} = (2M\ell^{2} - r^{3})/(r\ell^{2})$ is nonnegative on $r > 0$ exactly for $r^{3} \le 2M\ell^{2} = r_{*}^{3}$, and $Q(K) = 1 - \beta^{2}$ in this gauge (the form of the gauge and of $K$ are those of Theorem \claimref{synthesis.sds-clausius-fixes-area-constant}, with the sign of the $r^{2}/\ell^{2}$ term reversed), so $Q(K) = f$ and $Q(K) = 1$ at $r_{*}$. $f'(r) = 2M/r^{2} + 2r/\ell^{2} > 0$, $f \to -\infty$ as $r \to 0$ and $f \to +\infty$ as $r \to \infty$, so there is exactly one positive zero $r_{h}$. At it $\beta^{2} = 1 - f = 1 > 0$, so $r_{h} < r_{*}$ since $\beta^{2}$ is decreasing; and $2M/r_{h} = 1 + r_{h}^{2}/\ell^{2} > 1$ gives $r_{h} < 2M$.

\emph{(i).} $f(r_{h}) = 0$ is $2M = r_{h}(1 + r_{h}^{2}/\ell^{2})$, and substituting it in $f'$ gives $f'(r_{h}) = r_{h}^{-1}(1 + 3r_{h}^{2}/\ell^{2})$. Along the family at fixed $\ell$, $r_{h}$ is an increasing function of $M$ with $dM/dr_{h} = \tfrac12(1 + 3x^{2})$, and by (T), $T\,d(\pi r_{h}^{2}) = \tfrac{1}{4\pi}r_{h}^{-1}(1 + 3x^{2})\,2\pi r_{h}\,dr_{h} = dM$. So $S = \pi r_{h}^{2}$ satisfies (Cl), and since (Cl) fixes $dS/dr_{h}$ on the connected interval $M > 0$, every solution is $\pi r_{h}^{2} + c$. As $M \to 0$, $r_{h} < 2M \to 0$, so $S \to c$. At $\ell \to \infty$ this is the computation of Theorem \claimref{synthesis.quarter-from-temperature}.

\emph{(ii).} $4\pi\ell T = \tau(x) := (1 + 3x^{2})/x$, with $\tau'(x) = (3x^{2} - 1)/x^{2}$ and $\tau^{2} - 12 = (1 - 3x^{2})^{2}/x^{2} \ge 0$. So $\tau$ decreases from $+\infty$ on $0 < x < 1/\sqrt{3}$ to its minimum $2\sqrt{3}$ and increases to $+\infty$ beyond: each value above the minimum is taken exactly twice, once on each side. Since $dM/dx > 0$, the sign of $dM/dT$ is the sign of $\tau'$.

\emph{(iii).} $F = M - TS = \tfrac12 r_{h}(1 + x^{2}) - \tfrac{1}{4}r_{h}(1 + 3x^{2}) = \tfrac14 r_{h}(1 - x^{2})$. On the small branch $x^{2} < \tfrac13 < 1$, so $F > 0$; on the large branch $F < 0$ exactly for $x > 1$. $\tau(1) = 4$ gives $T_{\mathrm{HP}} = 1/(\pi\ell)$, above $T_{\min}$ by the factor $4/2\sqrt{3}$. Under (FE) the equilibrium field at temperature $T$ is the one of least $F$ among the horizonless field ($F = 0$) and the black holes at that temperature: for $T < T_{\mathrm{HP}}$ every black hole has $F > 0$ (the large one because $\tau$ is increasing there, so $T < T_{\mathrm{HP}}$ is $x < 1$), and for $T > T_{\mathrm{HP}}$ the large one has $F < 0$. At $x = 1$, $M = \ell$ and $S = \pi\ell^{2}$.

\emph{(iv).} The zero of $f$ gives $\mu = r_{h}^{d-3}(1 + x^{2})$ and $f'(r_{h}) = (d-3)\mu/r_{h}^{d-2} + 2r_{h}/\ell^{2} = r_{h}^{-1}\bigl((d-3) + (d-1)x^{2}\bigr)$. With $dS/dr_{h} = \tfrac14(d-2)\Omega_{d-2}r_{h}^{d-3}$, $T\,dS = \tfrac{(d-2)\Omega_{d-2}}{16\pi}\bigl((d-3)r_{h}^{d-4} + (d-1)r_{h}^{d-2}/\ell^{2}\bigr)dr_{h}$, whose integral from $0$ is the stated $M$. Then $16\pi F/\Omega_{d-2}r_{h}^{d-3} = (d-2)(1 + x^{2}) - (d-3) - (d-1)x^{2} = 1 - x^{2}$. The minimum of $\bigl((d-3) + (d-1)x^{2}\bigr)/x$ is at $x^{2} = (d-3)/(d-1)$, with value $2\sqrt{(d-1)(d-3)}$, below the zero of $F$ at $x = 1$, where $4\pi\ell T = 2(d-2)$; the argument of (ii) and (iii) applies unchanged.

\emph{(v).} By Theorem \claimref{synthesis.euler-term-horizon-entropy-shift}(ii), which holds on both branches of $\Lambda$ and rests on Fact \claimref{synthesis.imported-wald-noether-charge-entropy}, under (W) a spherical cross-section has $S = A/4 - 3\pi\chi/2\Lambda$ with $\chi = 2$, that is $c = \pi\ell^{2}$, a constant of the kind allowed by (i). Then $F_{W} = \tfrac14 r_{h}(1 - x^{2}) - \tfrac14\ell^{2}r_{h}^{-1}(1 + 3x^{2}) = -\tfrac{\ell}{4x}\bigl(x^{4} + 2x^{2} + 1\bigr)$, negative for every $x > 0$. So under (FE) each black hole has lower free energy than the horizonless field, on both branches, and $F_{W}$ has no zero. At equal temperature the two radii satisfy $x_{s}x_{l} = \tfrac13$ (the two roots of $3x^{2} - \tau x + 1 = 0$), and with $g(x) = (1 + x^{2})^{2}/x$ one has $g(x_{l}) - g(x_{s}) = (3x_{l}^{2} - 1)^{3}/27x_{l}^{3} > 0$, so the large black hole has the lower $F_{W}$. The constant is positive, so $S \ge 0$ places no restriction on it; $S \to \pi\ell^{2} \ne 0$ as $M \to 0$ violates (Z).

\emph{Orientation, not used.} With $\beta = 1/T$, the script confirms $\beta F_{W} = -\tfrac{\ell^{2}}{64\pi}\,\beta\cdot 4\pi\!\int_{r_{h}}^{\infty}\!\tfrac{48M^{2}}{r^{6}}\,r^{2}dr$ for every sampled $x$: the Wald free energy is, up to that factor, the integral of the square of the Weyl part $48M^{2}/r^{6}$ of the field strength over the Euclidean exterior, which is what an action that is a square of the Cartan curvature evaluates to on a field whose Cartan curvature is its Weyl part. The coefficient $\tfrac{1}{64\pi}$ is matched, not derived, here.

The script checks exactly, at six rational horizon radii and four values of $\ell$: the zero of $f$, its monotonicity at sample points, $\beta^{2} = 1$ at the horizon locus and $r_{h}^{3} < r_{*}^{3}$, the temperature and the Clausius identity as rational functions; the sign of $\tau'$ on each branch, the identity $\tau^{2} - 12 = (1 - 3x^{2})^{2}/x^{2}$, the free energy and its sign and zero; for $d = 4, \dots, 7$ the Clausius energy, the free energy and the values at $x = 1$ and at the minimum; the Wald constant from $\chi = 2$ and $\Lambda = -3/\ell^{2}$, the closed form and sign of $F_{W}$, and the orientation identity.
